% !PRINTED 557
\weberpage{557}

%<<CONT
(10) mit der Permutationsgruppe $C$ isomorph ist. Also ist $C$ die \gesperrt{Galois'sche} Gruppe von $X = 0$.

Nach \S.\,163 lässt sich die Gleichung $X = 0$ auf eine Reihe von cyklischen Gleichungen niedrigeren Grades zurückführen, wenn man Functionen aufsucht, die zu den Theilern der Gruppe $C$ gehören. Die Theiler von $C$ sind aber, wenn $n - 1$ in zwei Factoren $e, f$ zerlegt, also
\[
n - 1 = ef
\]
gesetzt wird, die cyklischen Gruppen

\weq{11}{C_e = 1, \pi^e, \pi^{2e} \ldots \pi^{(f-1)e},}

und die Permutation $\pi^e$ ist darin gleichbedeutend mit der Substitution $(r, r_e)$.

Um zu einer zu $C_e$ gehörigen Function der Wurzeln von möglichst einfacher Form zu gelangen, betrachten wir nach \gesperrt{Gauss} die $f$-\gesperrt{gliedrigen Perioden}

\weq{12}{\begin{aligned}
\eta &= r + r_e + r_{2e} + \cdots + r_{(f-1)e} \\
\eta_1 &= r_1 + r_{e+1} + r_{2e+1} + \cdots + r_{(f-1)e+1} \\
&\cdots\cdots\cdots\cdots\cdots\cdots\cdots\cdots \\
\eta_{e-1} &= r_{e-1} + r_{2e-1} + r_{3e-1} + \cdots + r_{n-2},
\end{aligned}}

die aus der ersten von ihnen durch die Substitutionen $(r, r), (r, r_1), (r, r_2) \ldots (r, r_{e-1})$ hervorgehen. Wir bezeichnen die Grössen (12) auch als ein System \gesperrt{conjugirter Perioden}.

Jede dieser Functionen bleibt durch die Substitution $(r, r_e)$ ungeändert, und sie sind alle von einander verschieden, da sich eine Gleichung $\eta_h - \eta_k = 0$, wenn $h$ von $k$ verschieden ist, in der Form
\[
a r + a_1 r_1 + \cdots + a_{n-2} r_{n-2} = 0
\]
schreiben liesse, wo die $a, a_1 \ldots a_{n-2}$ nicht alle zugleich verschwindende ganze Zahlen sind ($+ 1$, $0$ oder $- 1$), was nach dem Theorem I nicht möglich ist. Jede der Perioden (12) ist also eine zu der Gruppe $C_e$ \gesperrt{gehörige} Function, und diese Grössen sind daher die Wurzeln einer irreducibeln ganzzahligen Gleichung $e^{\mathrm{ten}}$ Grades. Adjungirt man eine dieser Grössen, so hängt die Bestimmung von $r$ noch von einer Gleichung vom Grade $f$ ab.

Nach den allgemeinen Sätzen des vierzehnten Abschnittes kann jede Function der $r$, die die Substitutionen der Gruppe $C_e$ gestattet, rational durch jede der Grössen $\eta$ dargestellt werden. Hier können wir aber noch den folgenden Satz aufstellen:
%CONT>>
